\documentclass[a4paper]{article}     
\usepackage{amsfonts}
\usepackage{amsmath}
\usepackage{amssymb}
\usepackage{graphicx}
\usepackage{color}

\newcommand{\Bgsin}{\operatorname{Bgsin}}
\newcommand{\Bgtan}{\operatorname{Bgtan}}
\newcommand{\Bgcos}{\operatorname{Bgcos}}
\newcommand{\Bgcot}{\operatorname{Bgcot}}

\begin{document}

\noindent \large Auteur: Thomas Olszowka \\
\noindent \large Studierichting: Informatica\\
\noindent \large Oefeningenreeks 2B nr. 1.3\\

\medskip

\normalsize

$f(x) = \sin(\Bgcot x)$ \\

neem $y = \Bgcot x$ \\

dan is $\cot y = x$ en $y \in \left]0, \pi\right[ $\\

we kunnen hier de identiteit $\sin^2y + \cos^2y  = 1$ gebruiken om de cotangens te vinden: \\

we weten dat $\cot y = \dfrac{\cos y}{\sin y}$ \\

dan is $\cot^2 y = \dfrac{\cos^2 y}{\sin^2 y}$ \\

$\sin^2y + \cos^2y  = 1$\\

$\dfrac{\sin^2 y}{\sin^2 y} + \dfrac{\cos^2 y}{\sin^2 y} = \dfrac{1}{\sin^2 y}$ \\

$1 + \cot^2 y = \dfrac{1}{\sin^2 y}$ \\

$\dfrac{1}{\sin^2 y} = 1 + \cot^2 y$\\

$\dfrac{1}{\sin^2 y}= 1 + x^2$  (als $\cot y = x$ is $\cot^2 y = x^2$) \\

$\sin^2 y = \dfrac{1}{1 + x^2}$\\

$\sin y = \pm \dfrac{\sqrt{1}}{\sqrt{1 + x^2}}$ \\

$= +\dfrac{1}{\sqrt{1 + x^2}}$ want $y \in \left]0, \pi\right[$\\

dus is $\sin(\Bgcot x) = \dfrac{1}{\sqrt{1 + x^2}}$

\begin{thebibliography}{999}
\end{thebibliography}
\end{document} 